% !TeX encoding = UTF-8

% 载入 SJTUThesis 模版
\documentclass[type=doctor]{sjtuthesis}
% 选项
%   type=[doctor|master|bachelor|course],     % 可选（默认：doctor），论文类型
%   zihao=[-4|5],                             % 可选（研究生默认：-4，本科默认：5），正文字号大小
%   lang=[zh|en],                             % 可选（默认：zh），论文的主要语言
%   review,                                   % 可选（默认：关闭），盲审模式
%   [twoside|oneside]                         % 可选（默认：twoside），单双页模式

% 论文基本配置，加载宏包等全局配置

% !TEX root = ./main.tex

\sjtusetup{
  %
  %******************************
  % 注意：
  %   1. 配置里面不要出现空行
  %   2. 不需要的配置信息可以删除
  %******************************
  %
  % 信息录入
  %
  info = {%
    %
    % 标题
    %
    title           = {基于深度学习的大规模图像检索},
    title*          = {Large-scale Image Retrieval Based on Deep Learning},
    %
    % 标题页标题
    %   可使用“\\”命令手动控制换行
    %
    % display-title   = {上海交通大学学位论文\\ \LaTeX{} 模板示例文档},
    % display-title*  = {A Sample Document \\ for \LaTeX-based SJTU Thesis Template},
    %
    % 页眉标题
    %
    % running-title   = {示例文档},
    % running-title*  = {Sample Document},
    %
    % 关键词
    %
    keywords        = {深度学习, 深度哈希, 乘积量化，行人重识别},
    keywords*       = {SJTU, master thesis, XeTeX/LaTeX template},
    %
    % 姓名
    %
    author          = {陈勇彪},
    author*         = {Yongbiao Chen},
    %
    % 指导教师
    %
    supervisor      = {戚正伟},
    supervisor*     = {Prof. Zhengwei Qi},
    %
    % 副指导教师
    %
    % assoc-supervisor  = {某某教授},
    % assoc-supervisor* = {Prof. Uom Uom},
    %
    % 学号
    %
    id              = {016037910002},
    %
    % 学位
    %   本科生不需要填写
    %
    degree          = {工学博士},
    degree*         = {Ph.D. of Engineering},
    %
    % 专业
    %
    major           = {软件工程},
    major*          = {Software Engineering},
    %
    % 所属院系
    %
    department      = {电子信息与电气工程学院, 软件学院},
    department*     = {School of Software, SEIEE},
    %
    % 课程名称
    %   仅课程论文适用
    %
    course          = {某某课程},
    %
    % 答辩日期
    %   使用 ISO 格式 (yyyy-mm-dd)；默认为当前时间
    %
    % date            = {2014-12-17},
    %
    % 资助基金
    %
    % fund  = {
    %           {国家 973 项目 (No. 2025CB000000)},
    %           {国家自然科学基金 (No. 81120250000)},
    %         },
    % fund* = {
    %           {National Basic Research Program of China (Grant No. 2025CB000000)},
    %           {National Natural Science Foundation of China (Grant No. 81120250000)},
    %         },
  },
  %
  % 风格设置
  %
  style = {%
    %
    % 本科论文页眉 logo 颜色 (red/blue/black)
    %
    % header-logo-color = black,
  },
  %
  % 名称设置
  %
  name = {
    % bib               = {References},
    % acknowledgements  = {谢\hspace{\ccwd}辞},
    % publications      = {攻读学位期间完成的论文},
  },
}

% 使用 BibLaTeX 处理参考文献
%   biblatex-gb7714-2015 常用选项
%     gbnamefmt=lowercase     姓名大小写由输入信息确定
%     gbpub=false             禁用出版信息缺失处理
\usepackage[backend=biber,style=gb7714-2015]{biblatex}
% 文献表字体
% \renewcommand{\bibfont}{\zihao{-5}}
% 文献表条目间的间距
\setlength{\bibitemsep}{0pt}
% 导入参考文献数据库
\addbibresource{bibdata/thesis.bib}

% 脚注格式
\usepackage[perpage,bottom,hang]{footmisc}

% 定义图片文件目录与扩展名
\graphicspath{{figures/}}
\DeclareGraphicsExtensions{.pdf,.eps,.png,.jpg,.jpeg}

% 确定浮动对象的位置，可以使用 [H]，强制将浮动对象放到这里（可能效果很差）
% \usepackage{float}

% 固定宽度的表格
% \usepackage{tabularx}

% 使用三线表：toprule，midrule，bottomrule。
\usepackage{booktabs}

% 表格中支持跨行
\usepackage{multirow}

\usepackage{amsmath}

% 表格中数字按小数点对齐
\usepackage{dcolumn}
\newcolumntype{d}[1]{D{.}{.}{#1}}

% 使用长表格
\usepackage{longtable}

% 附带脚注的表格
\usepackage{threeparttable}

% 附带脚注的长表格
\usepackage{threeparttablex}

% 算法环境宏包
\usepackage[ruled,vlined,linesnumbered]{algorithm2e}
% \usepackage{algorithm, algorithmicx, algpseudocode}


% 代码环境宏包
\usepackage{listings}
\lstnewenvironment{codeblock}[1][]%
  {\lstset{style=lstStyleCode,#1}}{}

% 物理科学和技术中使用的数学符号，定义了 \qty 命令，与 siunitx 3.0 有冲突
% \usepackage{physics}

% 直立体数学符号
\providecommand{\dd}{\mathop{}\!\mathrm{d}}
\providecommand{\ee}{\mathrm{e}}
\providecommand{\ii}{\mathrm{i}}
\providecommand{\jj}{\mathrm{j}}

% 国际单位制宏包
\usepackage{siunitx}[=v2]

% 定理环境宏包
\usepackage{ntheorem}
% \usepackage{amsthm}



% 绘图宏包
\usepackage{tikz}
\usetikzlibrary{shapes.geometric, arrows}

% 一些文档中用到的 logo
\usepackage{hologo}
\providecommand{\XeTeX}{\hologo{XeTeX}}
\providecommand{\BibLaTeX}{\textsc{Bib}\LaTeX}

% 借用 ltxdoc 里面的几个命令方便写文档
\DeclareRobustCommand\cs[1]{\texttt{\char`\\#1}}
\providecommand\pkg[1]{{\sffamily#1}}

% 自定义命令

% E-mail
\newcommand{\email}[1]{\href{mailto:#1}{\texttt{#1}}}

% hyperref 宏包在最后调用
\usepackage{hyperref}

% 自动引用题注更正为中文
\def\equationautorefname{式}
\def\footnoteautorefname{脚注}
\def\itemautorefname{项}
\def\figureautorefname{图}
\def\tableautorefname{表}
\def\partautorefname{篇}
\def\appendixautorefname{附录}
\def\chapterautorefname{章}
\def\sectionautorefname{节}
\def\subsectionautorefname{小节}
\def\subsubsectionautorefname{小节}
\def\paragraphautorefname{段落}
\def\subparagraphautorefname{子段落}
\def\FancyVerbLineautorefname{行}
\def\theoremautorefname{定理}


\begin{document}

%TC:ignore

% 标题页
\maketitle

% 原创性声明及使用授权书
\copyrightpage
% 插入外置原创性声明及使用授权书
% \copyrightpage[scans/sample-copyright-old.pdf]

% 前置部分
\frontmatter

% 摘要
% !TEX root = ../main.tex

\begin{abstract}
互联网技术，多媒体技术以及数字成像技术的蓬勃发展导致了网络上图像内容数据的爆炸性增长， 促使了图像检索技术的出现以及日渐成熟。
图像检索旨在从一个大规模的图像数据库中检索出和查询的图片对应的图像数据，是计算机视觉和多媒体检索领域一个核心的研究问题，因为其具有很高的科研价值以及广泛的应用，近些年在研究社区中获得了广泛的关注。
一个典型的图像检索系统包含两个阶段: (1) 表征学习: 学习一个算法提取图片物体的特征，将高维的图片数据映射到地维的特征向量空间。 (2) 度量学习: 使得第一步中学习到的特征向量具备高维图像数据的语义一致性。语义相似的图片的特征向量距离尽量小，相反，语义不同的的图像的特征向量距离尽量大。传统的图像检索使用实值的特征向量来表征图片, 虽然可以取得较高的准确度，但是由于内存占用量高以及检索速度不高, 导致其可扩展性差, 难以应用在对检索速度要求较高的大规模的图像检索场景。 \par
深度哈希是一种利用深度神经网络来进行特征提取，将高维图片数据映射到低维二进制向量空间的技术，因为其具有内存占用量低，并且检索速度快等优点而被广泛应用在图像检索领域来提高图像检索系统的可扩展性。本文致力于探究深度哈希算法在图像检索领域的应用，提高检索系统的精度和可扩展性。通过探究新型神经网络结构在检索技术中的应用来提高表征能力，以及探究新型度量学习在深度哈希技术中的应用， 本文提出了若干应用于图像检索领域的新型算法。 本文的主要创新点和贡献主要在以下几点:

\begin{enumerate}
  \item 本文对图像检索的一个典型应用-行人重识别问题进行了深入的探讨。传统的行人重识别问题针对单模态的数据，也就是检索的图片和数据库中的图片来自同一模态。而现实中的监控摄像头一般能同时在日间和夜间工作捕捉不同模态的行人照片，这使得跨模态行人重识别成为计算机视觉领域一个热门的研究课题。由于模态鸿沟的存在（Modality Gap）， 跨模态行人重识别比单模态行人重识别更具有挑战性。传统基于双流卷积神经网络的算法，很难学习到模态恒定（Modality-invariant）的特征。针对这一难点，本文创新性的基于自编码器（Autoencoder）设计了一个神经网络架构学习模态恒定以及外表恒定的特征。同时设计了适用于跨模态检索对其不同模态特征的损失函数，使得网络可以在动态创建的图片对上进行监督训练。我们在标准的跨模态重识别数据集上进行了训练以及测试，结果表明我们的算法可以取得优越的性能数据，证明了算法的有效性。
  \item 本文首次将深度哈希引入了大规模车辆检索的研究中。传统基于实值向量检索的车辆重识别算法虽然精度高，但是由于其存储效率低以及检索速度慢，无法适用于真实环境的大规模车辆检索。为了提高车辆检索的速度以及优化存储开销，本文第一个探索了深度哈希在大规模车辆检索中的应用。本文提出一个新型的离散哈希模块，通过同时使用传统的基于困难三元组的损失函数进行特征学习，以及离散哈希模块生成离散的汉明哈希码。为了优化离散的哈希架构，我们提出了一个交替优化算法来进行整个架构的优化。 本文在主流的车辆重识别数据集上进行了四种不同长度哈希码的精度测试，显著的超越了当前的普适的哈希算法。
  \item 为了进一步优化深度哈希算法的表征能力，我们第一次探索一个不基于卷积神经神经网络的深度哈希框架。视觉转换器（Vision Transformer）是一种基于自注意力机制（self-attention）的新型的计算机视觉基础网络模型。本文提出了第一个完全基于视觉转换器的深度哈希框架，通过设计一个孪生视觉转换器架构，以及一个新型的双流特征学习的视觉转换器模块来进行细粒度表征学习。 同时，我们采取成对的基于贝叶斯的学习框架进行度量学习。本文在三个标准化的图像检索基准数据集上进行测试，实验效果表明该方法可以大幅度提升检索的性能。
  \item 由于基于传统哈希编码的方法会带来较大的精度损失，本文探索一种基于乘积量化（Product Quantization）编码的深度哈希算法来提高检索性能。同时，本文提出第一个基于视觉转换器的乘积量化神经网络。为了进行细粒度的特征学习，考虑到视觉转换器的性能与采取的图片分块策略紧耦合的特点，本文设计了一个双支视觉转换器量化的支柱网络。同时，本文设计了一个基于排序损失的直接优化平均查准率 （Average Precision）的量化损失函数来进行度量学习。文章将乘积量化集成到端到端的神经网络优化中，取得了优异的性能表现。
\end{enumerate}




\end{abstract}

\begin{abstract*}
The remarkable advancements in internet , multimedia  and digital imaging technologies have prompted the explosive growth of multimedia data on the internet, which has spurred the advent and the rapid development of image retrieval technologies. Image retrieval is a key research topic in computer vision and multimedia retrieval which targets at retrieving the corresponding image in a large image datab ase given a specific query target. Substantial attention has been drawn to the research of image retrieval in the community given its notable research value and its widespread application in the industry. A common learning paradigm of image retrieval typically incoporates two phases. (1) representation learning: It aims to learn discriminative feature representation from the image data, mapping the high-dimensional image data into low-dimensional feature vector space.  (2) metric learning: it targets at preserving the semantic consistency of the images in the low-dimensional feature space. The projected feature vectors of the semantically-similar images should be pulled as close as possible while those of  dissimilar images be pushed far away. Conventional image retrieval methods generally adopt real-valued feature vectors which enjoy the merit of being sufficently accurate. Nontheless, their high storage costs and low retrieval speed have limited their applicability in a large-scale image retrieval setting where the retrieval speed is of critical importance. \par
Deep hashing is a pervasive technology which has been widespredly applied in image retrieval to promote the scalability of existing retrieval systems owing to its low memory usage and fast retrieval speed. Generally, it transforms high-dimensional image data into low-dimensional compact binary codes with deep neural networks. This dissertation focus on investigating the application of deep hashing technology in image retrieval to boost the accuracy and the scalability of existing retrieval systems.  By exploring novel deep learning architectures and metric learning algorithms, we propose several novel algorithms for large-scale image retrieval. The main novelties and contributions are stated as follows:
\begin{enumerate}
  \item We conduct in-depth research on a typical application of image retrieval - person re-identification. Conventional person re-identification focus on visible images captured by single-modality survelliance cameras. However, these visible light cameras fail to produce high-quality images under poor illumination conditions. Nowadays, the majority of the survelliance cameras are capable of automatically switching from visible to the infrared mode when the illumination condition is poor. This stimulates the researches on visible-infrared cross-modality person re-identification. The existence of the modality gap makes cross-modality person re-identification even more challenging. Previous methods with dual-stream convolutional neural networks generally fail to capture the modality-invariant features. To this end, we devise a novel modality and apperance invariant embedding learning framework for discriminative feature learning, which is equipped with retrofitted loss function for intra-modality and cross-modality retrieval. We conduct extensive experiments on standardized benchmark datasets which have evidenced the superiority of our proposed method.
  \item We make the very first attempt to adopt deep hashing learning to tackle large-scale vehicle retrieval problem. Conventional methods with real-valued feature vectors generally consume tremendous memory and computation, making them inapplicable in a real-world large-scale retrieval scenario. To enhance the retrieval speed and optimize the storage cost, we propose a deep hashing based vehicle re-identification framework.  Specifically, we propose a novel discrete hashing module, which could produce discrete hashing codes combined with the traditional feature learning module with hard-mining triplet loss. To optimize the overall architecture, we propose an alternating minimization method for similarity-preserving hashing codes learning.  Comprehensice experiments on standard vehicle re-identification datasets have demonstrated our effectiveness and efficiency in terms of four hashing code lengths.
  \item To further enhance the representation ability of current deep hashing algorithms, we explore the possibilty of devising a deep hashing framework without convolutional neural networks. Vision transformer is variant of transformer tailored for computer vision tasks, which is based on the self-attention mechanism. We introduce the first deep hashing framework based on pure vision transformer. We design a siamese vision transformer backbone for feature extraction and innovate a novel dual-stream feature learning block to learn discriminative global and local features. In addition, we adopt a bayesian learning scheme with a dynamic constructed similarity matrix for similarity-preserving learning. We performed experiments on three benchmark datasets and the results have evidenced our superiority against the state-of-the-art deep hashing methods.
  \item Traditional hashing encoding scheme will result in considerable performance declines. To mitigate this problem,  we investigate a deep hashing algorithm based on product quantization to boost the retrieval performance. We propose the fist product quantization network based on vision transformer. Motivated by the fact that the performance of vision transformer is closely-related with the patching strategy, we put forward the a dual-brnach vision transformer-based quantization network. In addition, we innovate a novel quantization loss, dubbed average precision quantization loss, which embedbed the asymmetric retrieval nature in product quantization-based methods into the metric learning process. The experiments on widely-studied benchmarks have demonstrated the effectiveness and superiority of our proposed quantization-based framework.
\end{enumerate}
\end{abstract*} 



% 目录
\tableofcontents
% 插图索引
\listoffigures*
% 表格索引
\listoftables*
% 算法索引
\listofalgorithms*

% 符号对照表
% !TEX root = ../main.tex

\begin{nomenclature*}
\label{chap:symb}

\sjtusetup{style={float-font=\normalsize}}

\begin{longtable}{rl}
  $\epsilon$  & 介电常数  \\  
  $\mu$       & 磁导率    \\
  $\epsilon$  & 介电常数  \\
  $\mu$       & 磁导率    \\
  $\epsilon$  & 介电常数  \\
  $\mu$       & 磁导率    \\
  $\epsilon$  & 介电常数  \\
  $\mu$       & 磁导率    \\
  $\epsilon$  & 介电常数  \\
  $\mu$       & 磁导率    \\
  $\epsilon$  & 介电常数  \\
  $\mu$       & 磁导率    \\
  $\epsilon$  & 介电常数  \\
  $\mu$       & 磁导率    \\
  $\epsilon$  & 介电常数  \\
  $\mu$       & 磁导率    \\
  $\epsilon$  & 介电常数  \\
  $\mu$       & 磁导率    \\
  $\epsilon$  & 介电常数  \\
  $\mu$       & 磁导率    \\
  $\epsilon$  & 介电常数  \\
  $\mu$       & 磁导率    \\
  $\epsilon$  & 介电常数  \\
  $\mu$       & 磁导率    \\
  $\epsilon$  & 介电常数  \\
  $\mu$       & 磁导率    \\
  $\epsilon$  & 介电常数  \\
  $\mu$       & 磁导率    \\
  $\epsilon$  & 介电常数  \\
  $\mu$       & 磁导率    \\
  $\epsilon$  & 介电常数  \\
  $\mu$       & 磁导率    \\
  $\epsilon$  & 介电常数  \\
  $\mu$       & 磁导率    \\
  $\epsilon$  & 介电常数  \\
  $\mu$       & 磁导率    \\
  $\epsilon$  & 介电常数  \\
  $\mu$       & 磁导率    \\
  $\epsilon$  & 介电常数  \\
  $\mu$       & 磁导率    \\
  $\epsilon$  & 介电常数  \\
  $\mu$       & 磁导率    \\
  $\epsilon$  & 介电常数  \\
  $\mu$       & 磁导率    \\
  $\epsilon$  & 介电常数  \\
  $\mu$       & 磁导率    \\
  $\epsilon$  & 介电常数  \\
  $\mu$       & 磁导率    \\
  $\epsilon$  & 介电常数  \\
  $\mu$       & 磁导率    \\
  $\epsilon$  & 介电常数  \\
  $\mu$       & 磁导率    \\
  $\epsilon$  & 介电常数  \\
  $\mu$       & 磁导率    \\
\end{longtable}

\end{nomenclature*}


%TC:endignore

% 主体部分
\mainmatter

% 正文内容
% !TEX root = ../main.tex

\chapter{绪论}
近些年来，由于互联网的普及，计算机技术的发展和各种社交网络的兴起，视频和图片数据以爆炸性的速度快速增长。如何对海量的图像数据进行处理以及分析促进了一系列的计算机视觉技术的飞速发展，图像检索技术是其中一个核心研究方向。早期基于文本的图像检索极大依赖于标注者的主观判断的标签，检索效率以及性能低下同时需要耗费大量的人力资源，难以应用于海量的互联网图像数据检索。基于内容的的图像检索（Content Based Image Retrieval, CBIR）通过算法自动提取图像的视觉特征来衡量图片之间的相似度，从而实现在数据库中进行图像检索的目的，极大的节省了人工标注耗费的人力、物力和财力。为了提高检索速度以及降低存储开销，哈希编码以及量化编码方法被广泛应用到了图像检索领域，成为解决超大规模图像检索问题的重要解决方案。 本章重点介绍大规模图像检索算法的研究背景以及意义，概要的介绍了图像检索算法发展历史, 详细介绍了基于深度哈希的大规模的图像检索算法。 本章最后介绍了这篇论文的章节安排以及概括介绍每个章节包含的主要内容。
\section{本研究的背景及意义}
在大数据、移动互联网、计算机技术、传感器技术、云计算等前沿技术的发展以及经济社会急速发展的需求下，人工智能（artificial intelligence）作为被视为驱动新一轮产业变革的技术受到了世界各国的广泛关注。 广义的人工智能\cite{russell2010artificial}
\section{国内外发展历史和研究现状}
\subsection{基于手工特征的图像检索}
\subsection{基于深度学习的主干网络发展}
\subsection{基于深度特征的图像检索}
\subsection{基于深度哈希的大规模检索}
\subsubsection{最近邻搜索}
\subsubsection{近似最近邻搜索}
\subsubsection{基于深度哈希编码搜索}
\subsubsection{基于深度学习的量化编码搜索}
\section{本文章节安排与研究内容}
% !TEX root = ../main.tex

\chapter{数学与引用文献的标注}

\section{数学}

\subsection{数字和单位}

宏包 \pkg{siunitx} 提供了更好的数字和单位支持：
\begin{itemize}
  \item \num{12345.67890}
  % For TeXLive 2021, siunitx >= 3.0
  % \item \complexnum{1+-2i}
  % For siunitx < 3.0
  % \item \num{1+-2i}
  \item \num{.3e45}
  % For TeXLive 2021, siunitx >= 3.0
  % \item \numproduct{1.654 x 2.34 x 3.430}
  % For siunitx < 3.0
  % \item \num{1.654 x 2.34 x 3.430}
  \item \si{kg.m.s^{-1}}
  \item \si{\micro\meter} $\si{\micro\meter}$
  \item \si{\ohm} $\si{\ohm}$
  \item \numlist{10;20}
  \item \numlist{10;20;30}
  \item \SIlist{0.13;0.67;0.80}{\milli\metre}
  \item \numrange{10}{20}
  \item \SIrange{10}{20}{\degreeCelsius}
\end{itemize}

\subsection{数学符号和公式}

按照国标 GB/T 3102.11—1993《物理科学和技术中使用的数学符号》，
微分符号 $\dd$ 应使用直立体。除此之外，数学常数也应使用直立体：
\begin{itemize}
  \item 微分符号 $\dd$：\cs{dd}
  \item 圆周率 $\uppi$：\cs{uppi}
  \item 自然对数的底 $\ee$：\cs{ee}
  \item 虚数单位 $\ii$, $\jj$：\cs{ii} \cs{jj}
\end{itemize}

公式应另起一行居中排版。公式后应注明编号，按章顺序编排，编号右端对齐。
\begin{equation}
  \ee^{\ii\uppi} + 1 = 0,
\end{equation}
\begin{equation}
  \frac{\dd^2 u}{\dd t^2} = \int f(x) \dd x.
\end{equation}

公式末尾是需要添加标点符号的，至于用逗号还是句号，取决于公式下面一句是接着公式说的，还是另起一句。
\begin{equation}
		\frac{2h}{\pi}\int_{0}^{\infty}\frac{\sin\left( \omega\delta \right)}{\omega}
		\cos\left( \omega x \right) \dd\omega = 
		\begin{cases}
				h, \ \left| x \right| < \delta, \\
				\frac{h}{2}, \ x = \pm \delta, \\
				0, \ \left| x \right| > \delta.
		\end{cases}
\end{equation}
公式较长时最好在等号“$=$”处转行。
\begin{align}
    & I (X_3; X_4) - I (X_3; X_4 \mid X_1) - I (X_3; X_4 \mid X_2) \nonumber \\
  = & [I (X_3; X_4) - I (X_3; X_4 \mid X_1)] - I (X_3; X_4 \mid \tilde{X}_2) \\
  = & I (X_1; X_3; X_4) - I (X_3; X_4 \mid \tilde{X}_2).
\end{align}

如果在等号处转行难以实现，也可在 $+$、$-$、$\times$、$\div$ 运算符号处转行，转行
时运算符号仅书写于转行式前，不重复书写。
\begin{multline}
  \frac{1}{2} \Delta (f_{ij} f^{ij}) =
    2 \left(\sum_{i<j} \chi_{ij}(\sigma_{i} - \sigma_{j})^{2}
    + f^{ij} \nabla_{j} \nabla_{i} (\Delta f) \right. \\
  \left. + \nabla_{k} f_{ij} \nabla^{k} f^{ij} +
    f^{ij} f^{k} \left[2\nabla_{i}R_{jk}
    - \nabla_{k} R_{ij} \right] \vphantom{\sum_{i<j}} \right).
\end{multline}

\subsection{定理环境}

示例文件中使用 \pkg{ntheorem} 宏包配置了定理、引理和证明等环境。用户也可以使用
\pkg{amsthm} 宏包。

这里举一个“定理”和“证明”的例子。
\begin{theorem}[留数定理]
\label{thm:res}
  假设 $U$ 是复平面上的一个单连通开子集，$a_1, \ldots, a_n$ 是复平面上有限个点，
  $f$ 是定义在 $U \backslash \{a_1, \ldots, a_n\}$ 上的全纯函数，如果 $\gamma$
  是一条把 $a_1, \ldots, a_n$ 包围起来的可求长曲线，但不经过任何一个 $a_k$，并且
  其起点与终点重合，那么：

  \begin{equation}
    \label{eq:res}
    \oint\limits_\gamma f(z)\, \dd z = 2\uppi \ii \sum_{k=1}^n \operatorname{I}(\gamma, a_k) \operatorname{Res}(f, a_k).
  \end{equation}

  如果 $\gamma$ 是若尔当曲线，那么 $\operatorname{I}(\gamma, a_k) = 1$，因此：

  \begin{equation}
    \label{eq:resthm}
    \oint\limits_\gamma f(z)\, \dd z = 2\uppi \ii \sum_{k=1}^n \operatorname{Res}(f, a_k).
  \end{equation}

  在这里，$\operatorname{Res}(f, a_k)$ 表示 $f$ 在点 $a_k$ 的留数，
  $\operatorname{I}(\gamma, a_k)$ 表示 $\gamma$ 关于点 $a_k$ 的卷绕数。卷绕数是
  一个整数，它描述了曲线 $\gamma$ 绕过点 $a_k$ 的次数。如果 $\gamma$ 依逆时针方
  向绕着 $a_k$ 移动，卷绕数就是一个正数，如果 $\gamma$ 根本不绕过 $a_k$，卷绕数
  就是零。

  定理~\ref{thm:res} 的证明。

  \begin{proof}
    首先，由……

    其次，……

    所以……
  \end{proof}
\end{theorem}

\section{引用文献的标注}

按照教务处的要求，参考文献外观应符合国标 GB/T 7714 的要求。模版使用 \BibLaTeX\
配合 \pkg{biblatex-gb7714-2015} 样式包%
\footnote{\url{https://www.ctan.org/pkg/biblatex-gb7714-2015}}
控制参考文献的输出样式，后端采用 \pkg{biber} 管理文献。

请注意 \pkg{biblatex-gb7714-2015} 宏包 2016 年 9 月才加入 CTAN，如果你使用的
\TeX\ 系统版本较旧，可能没有包含 \pkg{biblatex-gb7714-2015} 宏包，需要手动安装。
\BibLaTeX\ 与 \pkg{biblatex-gb7714-2015} 目前在活跃地更新，为避免一些兼容性问
题，推荐使用较新的版本。

正文中引用参考文献时，使用 \verb|\cite{key1,key2,key3...}| 可以产生“上标引用的参
考文献”，如 \cite{Yu2001,Cheng1999,LSC1957}。使用
\verb|\parencite{key1,key2,key3...}| 则可以产生水平引用的参考文献，例如
\parencite{Li1999,Jiang1989,Hopkinson1999}。请看下面的例子，将会穿插使用水平的和
上标的参考文献：普通图书\parencite{Yu2001,Jiang1998}，论文集、会议录
\cite{CSTAM1990}，科技报告\parencite{WHO1970}，学位论文\cite{Zhang1998}，专利文
献\parencite{Jiang1989,HBLZ2001}，专著中析出的文献\cite{Cheng1999,GBT2659}，期刊
中析出的文献\parencite{Li1999,Li2000}，报纸中析出的文献\cite{Ding2000}, 电子文献
\parencite{Jiang1999,Christine1998,Xiao2001}。

可以使用 \verb|\nocite{key1,key2,key3...}| 将参考文献条目加入到文献表中但不在正
文中引用。使用 \verb|\nocite{*}| 可以将参考文献数据库中的所有条目加入到文献表
中。
\nocite{Yang1999,Schinstock2000,Wen1990,GBT16159}

% !TeX root = ../main.tex

\chapter{浮动体}

\section{插图}

插图功能是利用 \TeX\ 的特定编译程序提供的机制实现的，不同的编译程序支持不同的图
形方式。有的同学可能听说“\LaTeX\ 只支持 EPS”，事实上这种说法是不准确的。\XeTeX
可以很方便地插入 EPS、PDF、PNG、JPEG 格式的图片。

一般图形都是处在浮动环境中。之所以称为浮动是指最终排版效果图形的位置不一定与源文
件中的位置对应，这也是刚使用 \LaTeX\ 同学可能遇到的问题。如果要强制固定浮动图形
的位置，请使用 \pkg{float} 宏包，它提供了 \texttt{[H]} 参数。

\subsection{单个图形}

图要有图题，研究生图题采用中英文对照，并置于图的编号之后，图的编号和图题应置于图
下方的居中位置。引用图应在图题右上角标出文献来源。当插图中组成部件由数字或字母等
编号表示时，可在插图下方添加图注进行说明，如图~\ref{fig:cn_100t} 所示。

\begin{figure}[!htp]
  \centering
  \includegraphics[width=4cm]{cn_100t.png} \\
    1.立柱 2.提升释放机构 3.标准冲击加速度计 \\
    4.导轨 5.重锤 6.被校力传感器 7.底座 \\
  \bicaption[出现在插图索引中]
    {单个图形示例\cite{He1999}。如果表格的标题很长，那么在表格索引中就会很不美观。可
      以在前面用中括号写一个简短的标题，这个标题会出现在索引中。}
    {Stay hungry, stay foolish.}
 \label{fig:cn_100t}
\end{figure}

\subsection{多个图形}

简单插入多个图形的例子如图~\ref{fig:SRR} 所示。这两个水平并列放置的子图共用一个
图形计数器，没有各自的子图题。

\begin{figure}[!htp]
  \centering
  \includegraphics[height=2cm]{sjtu-vi-badge-blue.pdf}
  \hspace{1cm}
  \includegraphics[height=2cm]{sjtu-vi-badge-blue.pdf}
  \bicaption{中文题图}{English caption}
  \label{fig:SRR}
\end{figure}

如果多个图形相互独立，并不共用一个图形计数器，那么用 \texttt{minipage} 或者
\texttt{parbox} 就可以，如图~\ref{fig:parallel1} 与图~\ref{fig:parallel2}。

\begin{figure}[!htp]
\begin{minipage}{0.48\textwidth}
  \centering
  \includegraphics[height=1.5cm]{sjtu-vi-name-blue.pdf}
  \caption{并排第一个图}
  \label{fig:parallel1}
\end{minipage}\hfill
\begin{minipage}{0.48\textwidth}
  \centering
  \includegraphics[height=1.5cm]{sjtu-vi-name-blue.pdf}
  \caption{并排第二个图}
  \label{fig:parallel2}
\end{minipage}
\end{figure}

如果要为共用一个计数器的多个子图添加子图题，建议使用较新的 \pkg{subcaption}宏
包，不建议使用 \pkg{subfigure} 或 \pkg{subfig} 等宏包。

推荐使用 \pkg{subcaption} 宏包的 \cs{subcaptionbox} 并排子图，子图题置于子图之
下，子图号用 a)、b) 等表示。也可以使用 \pkg{subcaption} 宏包的 \cs{subcaption}
（放在 minipage中，用法同 \cs{caption}）。

\pkg{subcaption} 宏包也提供了 \pkg{subfigure} 和 \pkg{subtable} 环境，如
图~\ref{fig:subfigure}。

\begin{figure}[!htp]
  \centering
  \begin{subfigure}{0.3\textwidth}
    \centering
    \includegraphics[height=2cm]{sjtu-vi-badge-blue.pdf}
    \caption{校徽}
  \end{subfigure}
  \hspace{1cm}
  \begin{subfigure}{0.4\textwidth}
    \centering
    \includegraphics[height=1.5cm]{sjtu-vi-name-blue.pdf}
    \caption{校名。注意这个图略矮些，subfigure 中同一行的子图在顶端对齐。}
  \end{subfigure}
  \caption{包含子图题的范例（使用 subfigure）}
  \label{fig:subfigure}
\end{figure}

搭配 \pkg{bicaption} 宏包时，可以启用 \cs{subcaptionbox} 和 \cs{subcaption} 的双
语变种 \cs{bisubcaptionbox} 和 \cs{bisubcaption}，如图~\ref{fig:bisubcaptionbox}
所示。

\begin{figure}[!hbtp]
  \centering
  \bisubcaptionbox{$R_3 = 1.5\text{mm}$ 时轴承的压力分布云图}%
                  {Pressure contour of bearing when $R_3 = 1.5\text{mm}$}%
                  [6.4cm]{\includegraphics[height=3cm]{example-image-a.pdf}}
  \hspace{1cm}
  \bisubcaptionbox{$R_3 = 2.5\text{mm}$ 时轴承的压力分布云图}%
                  {Pressure contour of bearing when $R_3 = 2.5\text{mm}$}%
                  [6.4cm]{\includegraphics[height=3cm]{example-image-b.pdf}}
  \bicaption{包含子图题的范例（使用 subcaptionbox）}
            {Example with subcaptionbox}
  \label{fig:bisubcaptionbox}
\end{figure}


\section{表格}

\subsection{基本表格}

编排表格应简单明了，表达一致，明晰易懂，表文呼应、内容一致。表题置于表上，研究生
学位论文可以用中、英文两种文字居中排写，中文在上，也可以只用中文。

表格的编排建议采用国际通行的三线表\footnote{三线表，以其形式简洁、功能分明、阅读
方便而在科技论文中被推荐使用。三线表通常只有 3 条线，即顶线、底线和栏目线，没有
竖线。}。三线表可以使用 \pkg{booktabs} 提供的 \cs{toprule}、\cs{midrule} 和
\cs{bottomrule}。它们与 \pkg{longtable} 能很好的配合使用。

\begin{table}[!hpt]
  \caption[一个颇为标准的三线表]{一个颇为标准的三线表\footnotemark}
  \label{tab:firstone}
  \centering
  \begin{tabular}{@{}llr@{}} \toprule
    \multicolumn{2}{c}{Item} \\ \cmidrule(r){1-2}
    Animal & Description & Price (\$)\\ \midrule
    Gnat  & per gram  & 13.65 \\
          & each      & 0.01 \\
    Gnu   & stuffed   & 92.50 \\
    Emu   & stuffed   & 33.33 \\
    Armadillo & frozen & 8.99 \\ \bottomrule
  \end{tabular}
\end{table}
\footnotetext{这个例子来自
  \href{https://mirrors.sjtug.sjtu.edu.cn/ctan/macros/latex/contrib/booktabs/booktabs.pdf}%
  {《Publication quality tables in LaTeX》}（\pkg{booktabs} 宏包的文档）。这也是
  一个在表格中使用脚注的例子，请留意与 \pkg{threeparttable} 实现的效果有何不
  同。}

\subsection{复杂表格}

我们经常会在表格下方标注数据来源，或者对表格里面的条目进行解释。可以用
\pkg{threeparttable} 实现带有脚注的表格，如表~\ref{tab:footnote}。

\begin{table}[!htpb]
  \bicaption{一个带有脚注的表格的例子}{A Table with footnotes}
  \label{tab:footnote}
  \centering
  \begin{threeparttable}[b]
     \begin{tabular}{ccd{4}cccc}
      \toprule
      \multirow{2}*{total} & \multicolumn{2}{c}{20\tnote{a}} & \multicolumn{2}{c}{40} & \multicolumn{2}{c}{60} \\
      \cmidrule(lr){2-3}\cmidrule(lr){4-5}\cmidrule(lr){6-7}
      & www & \multicolumn{1}{c}{k} & www & k & www & k \\ % 使用说明符 d 的列会自动进入数学模式，使用 \multicolumn 对文字表头做特殊处理
      \midrule
      & $\underset{(2.12)}{4.22}$ & 120.0140\tnote{b} & 333.15 & 0.0411 & 444.99 & 0.1387 \\
      & 168.6123 & 10.86 & 255.37 & 0.0353 & 376.14 & 0.1058 \\
      & 6.761    & 0.007 & 235.37 & 0.0267 & 348.66 & 0.1010 \\
      \bottomrule
    \end{tabular}
    \begin{tablenotes}
    \item [a] the first note.% or \item [a]
    \item [b] the second note.% or \item [b]
    \end{tablenotes}
  \end{threeparttable}
\end{table}

如某个表需要转页接排，可以用 \pkg{longtable} 实现。接排时表题省略，表头应重复书
写，并在右上方写“续表 xx”，如表~\ref{tab:performance}。

\begin{ThreePartTable}
  \begin{TableNotes}
    \item[a] 一个脚注
    \item[b] 另一个脚注
  \end{TableNotes}
  \begin{longtable}[c]{c*{6}{r}}
    \bicaption{实验数据}{Experimental data}
    \label{tab:performance} \\
    \toprule
    测试程序 & \multicolumn{1}{c}{正常运行} & \multicolumn{1}{c}{同步}
      & \multicolumn{1}{c}{检查点} & \multicolumn{1}{c}{卷回恢复}
      & \multicolumn{1}{c}{进程迁移} & \multicolumn{1}{c}{检查点} \\
    & \multicolumn{1}{c}{时间 (s)} & \multicolumn{1}{c}{时间 (s)}
      & \multicolumn{1}{c}{时间 (s)} & \multicolumn{1}{c}{时间 (s)}
      & \multicolumn{1}{c}{时间 (s)} &  文件（KB）\\
    \midrule
    \endfirsthead
    \multicolumn{7}{r}{续表~\thetable} \\
    \toprule
    测试程序 & \multicolumn{1}{c}{正常运行} & \multicolumn{1}{c}{同步}
      & \multicolumn{1}{c}{检查点} & \multicolumn{1}{c}{卷回恢复}
      & \multicolumn{1}{c}{进程迁移} & \multicolumn{1}{c}{检查点} \\
    & \multicolumn{1}{c}{时间 (s)} & \multicolumn{1}{c}{时间 (s)}
      & \multicolumn{1}{c}{时间 (s)} & \multicolumn{1}{c}{时间 (s)}
      & \multicolumn{1}{c}{时间 (s)}&  文件（KB）\\
    \midrule
    \endhead
    \hline
    \multicolumn{7}{r}{续下页}
    \endfoot
    \insertTableNotes
    \endlastfoot
    CG.A.2 & 23.05 & 0.002 & 0.116 & 0.035 & 0.589 & 32491 \\
    CG.A.4 & 15.06 & 0.003 & 0.067 & 0.021 & 0.351 & 18211 \\
    CG.A.8 & 13.38 & 0.004 & 0.072 & 0.023 & 0.210 & 9890 \\
    CG.B.2 & 867.45 & 0.002 & 0.864 & 0.232 & 3.256 & 228562 \\
    CG.B.4 & 501.61 & 0.003 & 0.438 & 0.136 & 2.075 & 123862 \\
    CG.B.8 & 384.65 & 0.004 & 0.457 & 0.108 & 1.235 & 63777 \\
    MG.A.2 & 112.27 & 0.002 & 0.846 & 0.237 & 3.930 & 236473 \\
    MG.A.4 & 59.84 & 0.003 & 0.442 & 0.128 & 2.070 & 123875 \\
    MG.A.8 & 31.38 & 0.003 & 0.476 & 0.114 & 1.041 & 60627 \\
    MG.B.2 & 526.28 & 0.002 & 0.821 & 0.238 & 4.176 & 236635 \\
    MG.B.4 & 280.11 & 0.003 & 0.432 & 0.130 & 1.706 & 123793 \\
    MG.B.8 & 148.29 & 0.003 & 0.442 & 0.116 & 0.893 & 60600 \\
    LU.A.2 & 2116.54 & 0.002 & 0.110 & 0.030 & 0.532 & 28754 \\
    LU.A.4 & 1102.50 & 0.002 & 0.069 & 0.017 & 0.255 & 14915 \\
    LU.A.8 & 574.47 & 0.003 & 0.067 & 0.016 & 0.192 & 8655 \\
    LU.B.2 & 9712.87 & 0.002 & 0.357 & 0.104 & 1.734 & 101975 \\
    LU.B.4 & 4757.80 & 0.003 & 0.190 & 0.056 & 0.808 & 53522 \\
    LU.B.8 & 2444.05 & 0.004 & 0.222 & 0.057 & 0.548 & 30134 \\
    EP.A.2 & 123.81 & 0.002 & 0.010 & 0.003 & 0.074 & 1834 \\
    EP.A.4 & 61.92 & 0.003 & 0.011 & 0.004 & 0.073 & 1743 \\
    EP.A.8 & 31.06 & 0.004 & 0.017 & 0.005 & 0.073 & 1661 \\
    EP.B.2 & 495.49 & 0.001 & 0.009 & 0.003 & 0.196 & 2011 \\
    EP.B.4 & 247.69 & 0.002 & 0.012 & 0.004 & 0.122 & 1663 \\
    EP.B.8 & 126.74 & 0.003 & 0.017 & 0.005 & 0.083 & 1656 \\
    SP.A.2 & 123.81 & 0.002 & 0.010 & 0.003 & 0.074 & 1854 \\
    SP.A.4 & 51.92 & 0.003 & 0.011 & 0.004 & 0.073 & 1543 \\
    SP.A.8 & 31.06 & 0.004 & 0.017 & 0.005 & 0.073 & 1671 \\
    SP.B.2 & 495.49 & 0.001 & 0.009 & 0.003 & 0.196 & 2411 \\
    SP.B.4 \tnote{a} & 247.69 & 0.002 & 0.014 & 0.006 & 0.152 & 2653 \\
    SP.B.8 \tnote{b} & 126.74 & 0.003 & 0.017 & 0.005 & 0.082 & 1755 \\
    \bottomrule
  \end{longtable}
\end{ThreePartTable}

\section{算法环境}

算法环境可以使用 \pkg{algorithms} 宏包或者较新的 \pkg{algorithm2e} 实现。
算法~\ref{algo:algorithm} 是一个使用 \pkg{algorithm2e} 的例子。关于排版算法环境
的具体方法，请阅读相关宏包的官方文档。

\begin{algorithm}[htb]
  \caption{算法示例}
  \label{algo:algorithm}
  \small
  \SetAlgoLined
  \KwData{this text}
  \KwResult{how to write algorithm with \LaTeXe }

  initialization\;
  \While{not at end of this document}{
    read current\;
    \eIf{understand}{
      go to next section\;
      current section becomes this one\;
    }{
      go back to the beginning of current section\;
    }
  }
\end{algorithm}

\section{代码环境}

我们可以在论文中插入算法，但是不建议插入大段的代码。如果确实需要插入代码，建议使
用 \pkg{listings} 宏包。

\begin{codeblock}[language=C]
#include <stdio.h>
#include <unistd.h>
#include <sys/types.h>
#include <sys/wait.h>

int main() {
  pid_t pid;

  switch ((pid = fork())) {
  case -1:
    printf("fork failed\n");
    break;
  case 0:
    /* child calls exec */
    execl("/bin/ls", "ls", "-l", (char*)0);
    printf("execl failed\n");
    break;
  default:
    /* parent uses wait to suspend execution until child finishes */
    wait((int*)0);
    printf("is completed\n");
    break;
  }

  return 0;
}
\end{codeblock}

\input{contents/summary}

%TC:ignore

% 参考文献
\printbibliography[heading=bibintoc]

% 附录
\appendix

% 附录中图表不加入索引
\captionsetup{list=no}

% 附录内容
\input{contents/app_maxwell_equations}
\input{contents/app_flow_chart}

% 结尾部分
\backmatter

% 用于盲审的论文需隐去致谢、发表论文、科研成果、简历

% 致谢
\input{contents/acknowledgements}

% 发表论文及科研成果
% 盲审论文中，发表论文及科研成果等仅以第几作者注明即可，不要出现作者或他人姓名
% !TEX root = ../main.tex

\begin{achievements}

\subsection*{学术论文}

\begin{bibliolist}{00}
  \item Chen, Yongbiao, Sheng Zhang, and Zhengwei Qi. "Maenet: Boosting feature representation for cross-modal person re-identification with pairwise supervision." Proceedings of the 2020 International Conference on Multimedia Retrieval. 2020.
  \item Chen, Yongbiao, et al. "Transhash: Transformer-based hamming hashing for efficient image retrieval." Proceedings of the 2022 International Conference on Multimedia Retrieval. 2022.
  \item Chen, Yongbiao, et al. "DVHN: A Deep Hashing Framework for Large-scale Vehicle Re-identification." arXiv preprint arXiv:2112.04937 (2021).
  \item Liu, Fangxin, et al. "Sstdp: Supervised spike timing dependent plasticity for efficient spiking neural network training." Frontiers in Neuroscience (2021): 1413.
  \item Chen, Yongbiao, et al. "Supervised Contrastive Vehicle Quantization for Efficient Vehicle Retrieval." Proceedings of the 2022 International Conference on Multimedia Retrieval. 2022.
  \item Liu, Fangxin, et al. "Bit-transformer: Transforming bit-level sparsity into higher preformance in reram-based accelerator." 2021 IEEE/ACM International Conference On Computer Aided Design (ICCAD). IEEE, 2021.
  \item Liu, Fangxin, et al. "Sato: spiking neural network acceleration via temporal-oriented dataflow and architecture." Proceedings of the 59th ACM/IEEE Design Automation Conference. 2022.
  \item Liu, Fangxin, et al. "Ebsp: evolving bit sparsity patterns for hardware-friendly inference of quantized deep neural networks." Proceedings of the 59th ACM/IEEE Design Automation Conference. 2022.
  \item Liu, Fangxin, et al. "SoBS-X: Squeeze-Out Bit Sparsity for ReRAM-Crossbar-Based Neural Network Accelerator." IEEE Transactions on Computer-Aided Design of Integrated Circuits and Systems 42.1 (2022): 204
  \item Liu, Fangxin, et al. "PIM-DH: ReRAM-based processing-in-memory architecture for deep hashing acceleration." Proceedings of the 59th ACM/IEEE Design Automation Conference. 2022.
\end{bibliolist}

% \begin{bibliolist*}{00}
%   \item 第一作者. 中文核心期刊论文, 2007.
%   \item 第一作者. EI 国际会议论文, 2006.
% \end{bibliolist*}

% \subsection*{专利}

% \begin{bibliolist}{00}
%   \item 第一发明人，“永动机”，专利申请号202510149890.0
% \end{bibliolist}

% \begin{bibliolist*}{00}
%   \item 第一发明人，“永动机”，专利申请号XXXXXXXXXXXX.X
% \end{bibliolist*}

\end{achievements}


% 简历
% !TEX root = ../main.tex

\begin{resume}
  \subsection*{基本情况}
    陈勇彪，1995 年 03 月生于 湖南娄底。

  \subsection*{教育背景}
  \begin{itemize}
    \item 2016 年 07 月至今，上海交通大学，博士研究生，软件工程 专业
    \item 2012 年 09 月至 2016 年 07 月，西北工业大学，本科，软件工程专业
  \end{itemize}

  \subsection*{研究兴趣}
    \LaTeX{} 机器学习， 计算机视觉， 大规模图像检索， 行人重识别

  \subsection*{联系方式}
  \begin{itemize}
    \item 地址： 上海市闵行区东川路 800 号，200240
    \item E-mail: \email{chenyongbiao0319@sjtu.edu.cn}
  \end{itemize}
\end{resume}


% 学士学位论文要求在最后有一个大摘要，单独编页码
\input{contents/digest}

%TC:endignore

\end{document}
